% Local charged-configuration estimate used in the trace response bound.
\subsection{The local charged-configuration estimate}
\label{sec:ex-lattice-trace-local-charge}

We isolate the finite-volume algebra that turns a charged configuration
into a bounded field and a priced response. This is the local step used
before the separation-limit passage. Its quantitative premises are stated
as component bounds, so the proof does not hide the estimate inside a
single named receipt.

For $t>e$, put
\[
 W_a(t)=\exp\bigl(a(\log\log t)^2\bigr),
 \qquad c(t)=t^{-30}.
\]
Then $W_a(t)W_b(t)=W_{a+b}(t)$, $W_{-1}(t)\to0$, and $c(t)\to0$ as
$t\to\infty$.

Let $I$ be a finite set of lattice directions, $Y$ a finite-dimensional
real normed coefficient space, and let $\ell:Y^I\to\mathbb R$ be a
continuous linear flux functional. Let $s:Y^I\to[0,\infty)$ be a size
function on coefficient arrays satisfying
$\|A\|_\infty\le s(A)$ for every coefficient array $A$; this domination is
the only property of $s$ used below. Write
$\sigma=\theta^{-2}$ with $\theta>0$ for the local normalization. The
nonlinear local coefficient map has already been evaluated in the arrays
$B',K',D'\in Y^I$; these are outputs of that local construction, not images
under the core linear projection. Let $X,W,E,L\in Y^I$. Assume the exact
charged-configuration decomposition
\begin{equation}
 X=W+B'-K'+D'+E+L,
 \label{eq:ex-lattice-charge-decomposition}
\end{equation}
and the following local coefficient and flux receipts at scale $t$:
\begin{align}
 s(X)&\le\theta^2 W_{1/6}(t),
 \label{eq:ex-lattice-charge-size-input}\\
 \ell(\sigma K')&=0,
 &\ell(\sigma B')&=h/12,
 &|h-m|&\le c(t),
 \label{eq:ex-lattice-charge-price-identities}\\
 s(W)&\le\theta^2 W_{-1}(t),
 &s(D')&\le\theta^2 W_{-1}(t),
 \label{eq:ex-lattice-charge-error-one}\\
 s(E)&\le\theta^2 W_{-1}(t),
 &s(L)&\le\theta^2 W_{-1}(t).
 \label{eq:ex-lattice-charge-error-two}
\end{align}
Here $h$ is the finite-volume charge extracted from the distinguished
component, and $m$ is its target price. These are coefficient-size receipts
after the local map has been evaluated. No norm contraction of the nonlinear
map is assumed.

\begin{proposition}[Local field and flux bounds]
\label{prop:ex-lattice-charge-bounds}
Under the displayed assumptions,
\[
 \|\sigma X\|_\infty\le W_{1/6}(t),
 \qquad
 \left|\ell(\sigma X)-\frac m{12}\right|
 \le4\|\ell\|W_{-1}(t)+\frac{c(t)}{12}.
\]
\end{proposition}
\begin{proof}
The field bound follows from the size domination and
\eqref{eq:ex-lattice-charge-size-input}:
\[
 \|\sigma X\|_\infty\le\sigma s(X)=\theta^{-2}s(X)\le W_{1/6}(t).
\]
For the price, apply $\ell$ to the decomposition in
\eqref{eq:ex-lattice-charge-decomposition}. The $K'$ term vanishes and the
$B'$ term is $h/12$. Therefore
\[
 \ell(\sigma X)-\frac m{12}
 =\ell(\sigma W)+\ell(\sigma D')
   +\ell(\sigma E)+\ell(\sigma L)+\frac{h-m}{12}.
\]
For any array $A$, continuity gives
$|\ell(\sigma A)|\le\|\ell\|\sigma\|A\|_\infty\le\|\ell\|\theta^{-2}s(A)$. The triangle
inequality and
\eqref{eq:ex-lattice-charge-price-identities}--
\eqref{eq:ex-lattice-charge-error-two} give the claimed bound. There are
four error terms, so the coefficient is exactly four; no direction-count
factor is introduced because $s$ bounds the full array norm.
\end{proof}

In the selected trace response, $m$ is the finite-volume payment of the
configuration under consideration; $m/12$ is not exactly the limiting flux.
The construction supplies in addition the accuracy bound
\[
 |m+12\kappa|\le t^{-100},
\]
where $\kappa=s_G>0$ is the slope constant \eqref{eq:shoot-slope}. Together with the proposition,
this gives
\[
 |\ell(\sigma X)+\kappa|
   \le4\|\ell\|W_{-1}(t)+\frac{t^{-30}+t^{-100}}{12},
\]
and the right side tends to zero as $t\to\infty$. The proposition and
this consequence concern the evaluated array $\sigma X$ of one
finite-volume configuration. The response whose size and flux are used in
Subsubsection~\ref{subsec:ex-lattice-trace-weighted-route} is instead the limit,
along the separation ultrafilter at a fixed level, of transported
finite-volume fields. Passing to it requires two further inputs, which do
not follow from the proposition: a transport bound of size $W_{-1/4}(t)$
for the difference between the transported field and the core projection
of $\sigma X$ (the core projection does not increase the norm and does not
change the flux), and the comparison of the arch scale with $t$. With
these, the finite-volume bounds take the form
\eqref{eq:ex-lattice-trace-finite-response-price}, continuity of the fixed
linear maps transfers them to the ultrafilter limit, and division of the
size bound by the arch scale gives a quantity tending to zero by
\eqref{eq:ex-lattice-trace-age-over-arch}. That passage is carried out,
conditionally on these inputs, in
Subsubsection~\ref{subsec:ex-lattice-trace-weighted-route}. The component
receipts themselves are separate finite-volume producer statements; this
proposition does not prove them from the charged-configuration hypotheses.

\subsubsection{Constructing the charged coefficient from its local pieces}
\label{subsec:trace-charge-construction}

We give the finite construction which precedes the proposition. This also
fixes which nonlinear projections occur in its six terms. All pieces are
on the same finite-volume representation and use its threshold $\theta$;
write $H=W_1(t)$ and put
\[
 c_0=\theta^2H^{-2},\qquad c_1=\theta^2H^{1/8},\qquad
 \epsilon_0=\theta^2t^{-1/100}.
\]
Let $\Pi_P f$ denote the weighted local coefficient of a regular observable
$f$ on a support $P$ containing the section footprint. Its readout constant
$r\ge1$ and residual constant $a\ge1$ are fixed at the stage. If the
section price is $p\ge0$, set $K=(1+p)r+a$.

First, the elementary local coefficient is an actual observable. Set
\[
 z=\Pi_P f,\qquad h=f-\operatorname{present}_\theta(z),\qquad
 Q_f=\operatorname{primary}_\theta(z)+{\bf1}_{\{P\ \mathrm{healthy}\}}h.
\]
Here the primary includes its own footprint mask. On a configuration
healthy on $P$, the two masks are one and $Q_f=f$. The residual has zero
local coefficient by construction. If the healthy derivative bound of $f$
is $u\ge0$, coefficient recovery and residual subtraction give
\[
 |z|_{\rm coeff}\le ru,\qquad
 |h|_{\rm healthy}\le au.
\]
Put $M_P(u)=e^{\alpha\operatorname{cov}(P)}u$, where $\alpha$ is the
support exponent. The coefficient with coordinate $z$, one regular
residual $h$ on $P$, and zero forced part has coefficient-plus-remainder
price at most
\[
 ru+M_P(au)\le (r+a)M_P(u)
\]
and pointwise exponential-entry bound at most
$((1+p)r+a)u\le K M_P(u)$. These inequalities use
$\operatorname{cov}(P)\ge0$ and the nonnegative support exponent
$\alpha$ (in the application $\alpha\ge4$), hence $M_P(u)\ge u$.
The construction is measurable, gauge invariant and depends only on the
links and healthy flags on $P$: these facts hold for the local projection,
its primary and the masked residual separately. Adding such coefficients
adds both prices and their exponential entries exactly. Thus this is a
construction of one coefficient, not a choice of its coordinate alone.

For each direction $b$, let $f_b^0,f_b^1,f_b^2,f_b^3$ be respectively the
local residual, the distinguished current term, the truncated loop term
and the Taylor remainder of the selected local decomposition. Their
healthy derivative bounds on the common patch $P$ are respectively
$c_0,c_1,c_1,c_0$. Let $g_{b,j}$, $j\in I_b$, be its finite additional
terms, supported on $P_{b,j}$ in the same enclosing patch, with healthy
bounds $u_{b,j}$ and
\begin{equation}
 \sum_{j\in I_b}M_{P_{b,j}}(u_{b,j})\le\epsilon_0.
 \label{eq:trace-charge-tail-price}
\end{equation}
Construct the elementary coefficient of
$f_b^0+f_b^1-f_b^2+f_b^3$, with derivative bound $2c_0+2c_1$, and add
those of the $g_{b,j}$. Denote the resulting coefficient by $X_b$.
The healthy observable and its coordinate are exactly
\begin{align*}
 Q_{X_b}&=f_b^0+f_b^1-f_b^2+f_b^3+\sum_jg_{b,j}
                      &&\text{on the common healthy patch},\\
 X_b.x&=\Pi_P f_b^0+\Pi_P f_b^1-\Pi_P f_b^2+\Pi_P f_b^3
                                      +\sum_j\Pi_{P_{b,j}}g_{b,j}.
\end{align*}
The second equality uses linearity of the local Taylor projection on
regular functions on the same patch; it makes no assertion of global
linearity for different nonlinear local charts.
Both prices are at most
\[
 K\{M_P(2c_0+2c_1)+\epsilon_0\}.
\]
The fixed patch cover bound gives $M_P(1)\le A$ with
$A=\exp(\alpha\,\mathrm{TileNo})$. For large $t$,
$\epsilon_0\le c_0\le c_1$ and therefore
\[
 K\{M_P(2c_0+2c_1)+\epsilon_0\}
 \le(4A+1)K\theta^2H^{1/8}\le\theta^2H^{1/6}.
\]
Indeed the last ratio is $(4A+1)K H^{-1/24}\to0$, while
$\epsilon_0/c_0=t^{-1/100}H^2\to0$. The sum of added coordinates has
size at most $(r+a)\epsilon_0\le K\epsilon_0\le c_0$.
Also $rc_0\le\theta^2H^{-1}$ eventually. Thus the local residual,
Taylor remainder, and added-coordinate terms have the prices asserted in
\eqref{eq:ex-lattice-charge-error-one}--
\eqref{eq:ex-lattice-charge-error-two}.
These are fixed-stage estimates. To use them along increasing stages,
choose the threshold separately at each stage, simultaneously for the
finitely many grades at most five and all its fixed readout constants.
Lemma~\ref{lem:ex-lattice-monotone-threshold} gives the monotone delay
envelope. The selected paid schedule meets these thresholds before the
period, volume and finite truncations are chosen; this is not an assertion
that unbounded stage constants are controlled by one fixed-stage tail.
The coefficient has zero forced part. Its density is consequently the
background density times its own exponential entry. The zero paid forced
response has attainment cost zero, hence at most $t$. This supplies the
atomic certificate used in the Coffer application, with charge
$\theta^2H^{1/6}$.

To obtain the six-term decomposition of the proposition, use the exact
loop splitting
\[
 f_b^2=k_b-d_b
\]
on this same patch, where $d_b$ has healthy derivative bound $c_0$.
Define
\[
 W_b=\Pi_P f_b^0,\quad B'_b=\Pi_P f_b^1,\quad
 K'_b=\Pi_P k_b,\quad D'_b=\Pi_P d_b,\quad
 E_b=\Pi_P f_b^3,\quad L_b=\sum_j\Pi_{P_{b,j}}g_{b,j}.
\]
Local projection of the splitting gives $X=W+B'-K'+D'+E+L$.
The $D'$ bound follows from $rc_0\le\theta^2H^{-1}$ as well. The
coefficient charge bounds $|X_b.x|_{\rm coeff}$ by
$\theta^2H^{1/6}$; taking the maximum over directions proves the full
array size estimate, rather than summing that estimate four times.

The two remaining flux statements have a different origin from these
size bounds. The local radial reading identifies the normalized flux of
$\Pi_P f$ with one twelfth of the cubic current probe of $f$. Thus the
loop cancellation and distinguished-current comparison needed here are
\begin{equation}
 \mathfrak H(k)=0,\qquad
 |\mathfrak H(f^1)-m|\le t^{-30},\qquad
 \ell(\theta^{-2}\Pi_P f)=\mathfrak H(f)/12.
 \label{eq:trace-charge-current-identities}
\end{equation}
They give exactly the two identities in
\eqref{eq:ex-lattice-charge-price-identities}, with
$h=\mathfrak H(f^1)$. In particular the cancellation concerns $k$, not
$f^2$: the small loop correction $d$ must remain among the four errors.
A bound on $f^2$ alone would not prove the flux identity.

The following subsections construct the remaining local prices and
identities on these same fields: the corrected-loop cancellation and
counterterm size, the actual nonlinear localization, the support-weighted
tail, and the fifth-order residual and finite-difference remainder.
They use the selected history's inherited action equation and local
analytic probes. The exact finite tensor corrections are displayed
below, so the trace and antisymmetric calibration identities are not
additional assumptions. Once these constructions are combined, the
coefficient and all its prices above are produced by the displayed
finite sum.

\subsubsection{Why the corrected loop has zero current price}
\label{subsec:trace-charge-loop-cancellation}

The cancellation in \eqref{eq:trace-charge-current-identities} can be
checked before estimating any norms. A quadratic loop word has two
plaquette legs, at offsets $o_{j,0},o_{j,1}$, and a color coefficient
$c_j$. All sums below can be taken separately for each pair of plaquette
orientations and each pair of color indices. Put
$\delta_j=o_{j,1}-o_{j,0}$. The corrected finite loop family is chosen with
\begin{equation}
 \sum_jc_j=0,\qquad \sum_jc_j\delta_j=0.
 \label{eq:trace-charge-loop-moments}
\end{equation}
These conditions have a concrete finite construction: subtract the local
zeroth-moment word and the first-moment difference words from the original
finite family, separately in each orientation/color component. A difference
word with offsets $0,e_a$, minus the word with offsets $0,0$, has zero
zeroth moment and first moment $e_a$. Taking its coefficient to be the
$a$th component of the original first moment, and the local word's
coefficient to be the original zeroth moment, cancels both moments
exactly. The original loop equals the corrected family minus its added
counterterms. The color coefficients of every added word are linear combinations of
the original invariant color coefficients, so they remain invariant.
These counterterms are the observable denoted $d$ above;
they cannot be discarded when estimating its size.

Let $F$ be the translated and signed-averaged sum of the corrected words.
At the identity configuration every curvature factor is zero, so
$F(0)=DF(0)=0$. Write $dA$ for the linear plaquette curl of a link
variation $A$. Differentiating the product of two curvature factors twice
therefore gives, before signed averaging,
\[
 D^2F(0)[A,B]=\sum_{j,z}c_j\bigl(
 (dA)(z+o_{j,0})(dB)(z+o_{j,1})
 +(dB)(z+o_{j,0})(dA)(z+o_{j,1})\bigr).
\]
The displayed products retain their prescribed color contractions.
Suppose $A$ has bounded support and $dB$ is affine on the support enlarged
by the word lengths. Translate the first leg to $z$. The inner sum of
the first term is a multiple of
\[
 \sum_jc_j\{b_0+L(z)+L(\delta_j)\}
 =(b_0+L(z))\sum_jc_j+L\Bigl(\sum_jc_j\delta_j\Bigr)=0.
\]
The other term has offsets $-\delta_j$ and vanishes for the same reason.
Nonwrapping of the enlarged patch justifies these ordinary integer
translations; it avoids replacing an affine function by a discontinuous
periodic representative. Signed averaging preserves the conclusion, since
it preserves support and the property of having affine curl. Thus
\begin{equation}
 D^2F(0)[A,B]=0
 \quad\text{if $A$ is supported in the patch and $dB$ is affine there.}
 \label{eq:trace-charge-loop-hessian}
\end{equation}

Now let $Y_b$ be the local gauge-covariant vector field of the chosen
current, vanishing at the identity, and let $k_b=DF\,Y_b$. The support
of $DY_b(0)A$ is contained in the same patch. Twice applying the product
rule, using $DF(0)=0$ and $Y_b(0)=0$, gives
\[
 D^2k_b(0)[A,B]
 =D^2F(0)[A,DY_b(0)B]+D^2F(0)[B,DY_b(0)A].
\]
For the polynomial link probes of degrees zero, one and two entering the
local cubic current test, their curls are affine. Each term consequently
vanishes by \eqref{eq:trace-charge-loop-hessian} and symmetry of the
Hessian. The constant-direction terms require a separate argument. The function
$k_b=DF\,Y_b$ is gauge invariant, since $F$ is gauge invariant and $Y_b$
is gauge covariant. Moreover $k_b(0)=Dk_b(0)=0$. Differentiating gauge
invariance once in a gauge parameter and once in an arbitrary variation
therefore says that $D^2k_b(0)$ annihilates every infinitesimal gauge
direction. On the nonwrapping patch a constant link variation is such a
direction: use the gauge parameter linear in the lattice coordinate and
extend it outside the patch. This kills the $(0,3)$ and $(3,0)$ terms
without claiming that a cubic probe has affine curl.
The radial current reading is a finite linear combination of these
Hessian probes: its cubic spatial coefficient is the polarization of the
pairs of spatial degrees $(0,3),(1,2),(2,1),(3,0)$, with the constant-curl
terms removed by the gauge identity. The middle two pairs are exactly the
quadratic link probes just treated. Hence every term of its current
reading is zero and $\mathfrak H(k)=0$.
This proof uses moment cancellation of the corrected loop, rather than
smallness or a limiting Ward identity.

\subsubsection{The exact reference-profile price}
\label{subsec:trace-reference-local-calibration}

The stress probe and its divergence are read in grades four and five,
respectively. Grade four is particularly useful here: its nonconstant
invariant labels are quadratic in curvature, so the local reading is
determined by its Hessian on constant-curvature radial jets. A translated
finite difference has zero such reading. Indeed translation leaves these
Hessian tests unchanged, and the difference of the two equal readings is
zero. This statement concerns the local grade-four reading, not the
observable itself.

To make the same-object identities explicit, let $\widetilde J$ be
the componentwise localized tensor of the selected field, and let $Q_i$
and $A_{abi}$ be its localized trace and antisymmetric auxiliary
currents. Write $\Delta_i$ for the unit translated difference. The
corrected local tensor $M$ keeps $\widetilde J_{ab}$ for $a<b$ and
$\widetilde J_{aa}$ for $a\ne0$, and is defined on the remaining entries by
\[
 M_{ba}=M_{ab}-\sum_i\Delta_i A_{abi}\quad(a<b),\qquad
 M_{00}=\sum_i\Delta_iQ_i-\sum_{i\ne0}M_{ii}.
\]
Consequently, before any projection or averaging,
\[
 \sum_iM_{ii}=\sum_i\Delta_i Q_i,\qquad
 M_{ab}-M_{ba}=\sum_i\Delta_i A_{abi}\quad(a<b).
\]
All these auxiliary fields and their translated differences use the same
finite nonwrapping support and the same healthy regular construction.
Thus linearity of the weighted local projection applies to these exact
identities. Applying the preceding zero-reading observation proves that
its grade-four reading is symmetric and trace free. Signed averaging
preserves both statements: each diagonal receives the square of the sign,
and permutations only reorder the trace sum. The grade-four classification
in Lemma~\ref{lem:ex-trace-free-stress-basis} therefore writes the averaged
reading uniquely as the three-parameter nonconstant stress array
$\mathsf S(a)$, $a\in\mathbb R^3$. Reading the local divergence commutes
with taking the grade-five derivative: this follows by differentiating the
finite radial Taylor polynomial and comparing each jet coefficient.
Consequently its vector reading is the same $Va$ used in
\eqref{eq:trace-weighted-calibration-price}; it is not independently
calibrated.

The trace and vector price in these coordinates are
\[
 \sum_b\mathsf S(a)_{bb}=4(a_0+4a_1)\mathcal O,\qquad
 \rho(Va)=-4(a_0+4a_1).
\]
Since $\mathcal O\ne0$ for a nontrivial simple gauge group, the zero
trace gives $a_0+4a_1=0$, hence $\rho(Va)=0$ exactly at finite volume.
For the two elementary correction tensors the local reading has the
fixed triples $a^0,a^1$, so their prices equal $\rho(Va^P)$ exactly.
This proves the exact price part of
\eqref{eq:trace-weighted-calibration-price} without a quantitative limit
for the stress coefficients.

For the qualitative stress calibration, let $J_{bc}$ be the unmodified
stress probe and $\widetilde J_{bc}$ its local version. On each fixed
constant-curvature test pair the selected local estimates give
\[
 \big|D^2J_{bc}(0)[u,v]-D^2(-\mathsf S_{\rm Maxwell})_{bc}(0)[u,v]\big|
 \le t^{-100},\qquad
 \big|D^2(J_{bc}-\widetilde J_{bc})(0)[u,v]\big|
 \le16t^{-30}.
\]
The factor $16$ is the product of the two factor-four link bounds for a
constant-curvature radial test. There are finitely many tests determining
the grade-four frame: evaluation on all radial jets is injective, and in
finite dimension a finite subset already has zero common kernel.
Their inverse on its image has finite operator norm. The triangle
inequality therefore gives convergence of the entire local coefficient
to the negative Maxwell coefficient. Applying the continuous inverse of
$a\mapsto\mathsf S(a)$ gives $a\to(-1,1/4,-1)$, uniformly on any sets on
which these two test bounds hold. This is precisely the qualitative
reference limit used in the weighted estimate; its weighted price uses
the exact zero above and does not multiply this qualitative error by
$A_{1/5}$.

\subsubsection{The localization error in the current comparison}
\label{subsec:trace-charge-localization}

The current comparison uses a shrinking analytic width, so its loss
cannot be treated as a fixed constant. At scale $t$, the width is
$a=(\log t)^{-1}$ and the localization radius is
$R=1000\lceil(\log t)^2\rceil$. Here is the nonlinear localization estimate for the actual localizer.
Write $F$ for the scalar reading of a gauge-invariant local field, $J_4F$
for its Taylor polynomial of degree four, $S_R$ for restriction of link
coordinates to distance at most $R$, and $W$ for the degree-four local
gauge-coordinate map. These readings are analytic at the zero field and
satisfy $F=F\circ\mathrm{exactP}_o$ as germs there: this is the local
gauge reconstruction identity of the constructed invariant field.
The Taylor polynomial of that reconstruction is $W$. Thus the defining
coordinate identity gives equality of four-jets
\[
 F-(J_4F\circ S_R)\circ W
       \equiv_4 (J_4F-J_4F\circ S_R)\circ W.
\]
This is an identity of jets, not an assertion of equality of the full
nonlinear fields. In a coordinate derivative indexed by links
$(e_1,\ldots,e_j)$, $j\le4$, the jet in parentheses is zero if
$\max_i d(o,e_i)\le R$, and otherwise is the original derivative of $F$.
Thus every retained term has at least one input outside the radius.

For completeness, the local coordinate map has derivative row sums,
through order four and on a radius-$g$ set, bounded by
$B(1+g)^L\exp(4(a/100)(g+2))$, with $B,L$ depending only on the gauge
group and derivative order. To see the volume independence of these bounds, a coordinate at distance
$g$ depends only on coordinates within radius $4(g+2)$. There are at most
$2(2\cdot4(g+2)+1)^4\dim(\mathrm{Loc})$ such coordinates. Its derivatives
through order four have a polynomial bound in $g$ by the finite product
rule for the local gauge-coordinate map. Multiplying that bound by the
fourth power of this coordinate count and by the polynomial spatial
weight gives $B(1+g)^L$; the exponential weight contributes the displayed
exponential. Thus the count concerns only this finite neighborhood. The chain rule is a sum over the
ordered partitions of $j\le4$. For each partition, there are at most
four such row factors. Using $g>R$, their product, together with the
output weight of width $a/100$, is bounded by the original input weight
of width $a$, times
\[
 B^4 e^{100}\operatorname{momentPrice}(4L,a/4)e^{-aR/4}.
\]
Indeed $(1+g)^{4L}\le\operatorname{momentPrice}(4L,a/4)e^{ag/4}$,
and the remaining exponent satisfies
$(a/100)g+16(a/100)(g+2)+ag/4
 \le ag+100-aR/4$ for $0<a\le1$.
In the following display $\|\cdot\|_{a,\le4}$ denotes the maximum of
the exponentially weighted coordinate-derivative sums \emph{at zero},
through order four, with the degree-four polynomial spatial weight.
It is not a supremum of derivatives over the healthy domain.
Summing over the finitely many ordered partitions gives a constant $D$
and an integer $P$, independent of the volume, such that
\begin{equation}
 \|F-F^R\|_{a/100,\,\le4}
 \le D\operatorname{momentPrice}(P,a/4)e^{-aR/4}
          \|F\|_{a,\,\le4}.
 \label{eq:trace-charge-nonlinear-localization}
\end{equation}
The scalar readings of the actual vector and tensor localizers are
exactly $(J_4F\circ S_R)\circ W$: this follows by substituting their
linkwise polynomial definition and contracting the link index. Their
zero and first derivatives vanish; thus the weighted estimate for
orders two through four supplies all the input derivatives in the
calculation. This proves the estimate for those localizers, without
promoting a scalar convolution estimate to a nonlinear one.

The selected input price is bounded by $C H^q$ for a fixed $q$.
The moment price above is a fixed polynomial in $a^{-1}$.
Since $aR/4\ge250\log t$, the right side of
\eqref{eq:trace-charge-nonlinear-localization} is $t^{-250+o(1)}$,
and is eventually at most $t^{-98}$. Only orders at most four are used
here; they suffice for the second-derivative current test and its one
divergence derivative.

To make the current test estimate explicit, write $Y_b$ for the original
local vector field and $Y_b^R$ for its localized field. The divergence
operator defining the current has one more derivative and a finite
weighted link sum. Applying the product rule and the preceding tail
bound, with the analytic width decreased to $a/100$, gives an error at
most $t^{-70}$ in the weighted derivative norm needed for this operator.
The spare exponent absorbs all fixed derivative orders and polynomial
support factors. The current probe is a fixed finite sum of spatial
moments of its second derivatives. If $c=a/(100\cdot5\cdot4\cdot5)$,
the sums of the moment weights of degrees two and three are bounded by
$C(c^{-1}+1)^m$ for a fixed $m$: sum the four-dimensional geometric
series and differentiate it at most three times. Consequently
\[
 |\mathfrak H(\operatorname{div}Y)
      -\mathfrak H(\operatorname{div}Y^R)|
 \le C(1+\log t)^m t^{-70}\le t^{-30}
\]
for all sufficiently large $t$. Subtracting the constant value of
$\operatorname{div}Y^R$ does not change this test, which uses its second
derivatives. This proves the distinguished-current error in
\eqref{eq:trace-charge-current-identities} for the localized field.
The same argument with the two constant-curvature directions gives the
$16t^{-30}$ local stress-probe error used above. The other stress-probe
error, $t^{-100}$, is the chosen accuracy of the original finite-volume
probe; the two bounds concern different comparisons.

\subsubsection{Paying the full reference observable}
\label{subsec:trace-reference-atomic-certificate}

The reference profiles require their full observable, including its value
on unhealthy configurations. They therefore use the following guarded
construction rather than the zero-forced coefficient above. Let $f$ be one
of the reference tensor or divergence probes on $P$, with global bound
$b=H^{1/50}$ and healthy derivative bound $c=\theta^2H^{1/8}$. These are
the local regularity prices for this profile. Set $z=\Pi_P f$ and let
$h=f-\operatorname{present}_\theta z$ be its zero-reading residual.
The primary's footprint mask differs from the mask on $P$; keep this
difference by defining
\[
 q=f-\operatorname{primary}_\theta z
                       -{\bf1}_{\{P\ \mathrm{healthy}\}}h.
\]
Then the equality $f=\operatorname{primary}_\theta z+
{\bf1}_{\{P\ \mathrm{healthy}\}}h+q$ holds on every configuration.
Here $a=1+(1+\mathrm{readable})r$ is the canonical atomic price,
so the local residual construction bounds $h$ by $ac$ on the healthy set
and bounds the footprint-mask discrepancy by $ac$ globally. Hence
$q$ is zero on the healthy set and $|q|\le b+ac$ everywhere. It is
measurable, invariant, supported on $P$ and depends only on its links and
flags, so $q$ is an admissible supported forced response. Its attainment
price is
\[
 e^{(\alpha+2)\operatorname{cov}(P)}(b+ac).
\]

The retained forced-activity expansion installs this supported response
in the same background density. To recall its quantitative content, mark
one copy of $q$ on $P$, expand the background activities and collect
each nonempty subset meeting $P$ into a supported forced slot, grouping
by its union with $P$. Untouched activities remain outside this slot.
There is no empty-slot contribution: $q$ vanishes when $P$ is healthy.
If the rooted
forced budget is $v$, their total absolute price is bounded by
\[
 M_P(b+ac)(e^{|P|v}-1).
\]
Indeed for $j$ selected factors the rooted product sum is at most
$(|P|v)^j/j!$; sum over $j\ge1$. This is the same supported forcing
expansion used by the atomic input of
Subsection~\ref{sec:ex-lattice-trace-coffer-supplier}; all terms retain the
same marked origin and its flags. Its exact product expansion makes the
resulting coefficient density equal to $\operatorname{eval}\,f$, not
just equal on the healthy set. For $|P|v\le1$,
$e^{|P|v}-1\le2|P|v$. In the remainder norm with tolerance $\sqrt v$
the resulting price is consequently at most
$2|P|\sqrt v\,M_P(b+ac)$. Together with the ordinary residual and the
coefficient this gives
\begin{equation}
 |z|_{\rm coeff}+\mathrm{rem}
 \le rc+M_P(ac)+2|P|\sqrt v\,M_P(b+ac).
 \label{eq:trace-reference-guarded-price}
\end{equation}
This explains both the reserved support exponent $\alpha+2$ in the
attainment price and the extra masked-remainder term; dropping either
would not give the reference atomic certificate.

Here the actual forced budget is
$v=\exp\{-\exp[(\log\log(t^{1/5}))^2]\}$,
$|P|=O((\log t)^8)$ and
$\operatorname{cov}(P)\le\mathrm{TileNo}$. The last term of
\eqref{eq:trace-reference-guarded-price} is smaller than every negative
power of $t$: take logarithms and use
$\exp[(\log\log(t^{1/5}))^2]/\log t\to\infty$. It is eventually at most $\theta^2$.
The other terms are at most $(r+Aa)\theta^2H^{1/8}$, so the whole price
is at most $\theta^2H^{1/6}$ after the stage delay has been paid.
Similarly
\[
 e^{(\alpha+2)\mathrm{TileNo}}(H^{1/50}+a\theta^2H^{1/8})\le t
\]
eventually. Thus each reference profile has the atomic charge and
attainment price used in \eqref{eq:trace-weighted-reference-bias}, with
its original full observable and the same weighted local coefficient.
Taking the maximum of the finitely many thresholds covers all three
profiles and their tensor and divergence entries.

\subsubsection{The weighted nonempty-support tail}
\label{subsec:trace-charge-support-tail}

The added supports in \eqref{eq:trace-charge-tail-price} are the
nonempty subsets of the activity increments in the finite shifted
product. Here is the estimate which keeps their support weights intact.
Fix a marked link, let $Q$ be the core support, and let $X_i$ be the
supports of its activity increments $f_i$, all containing the marked
site. Write $p_i$ for their derivative prices and $A$ for the price of
the unshifted local factor. If the derivative order is $k$, set
$c_k=2^k\mathrm{budC}(k)$. The product rule and the derivative series for
$e^{f_i}-1$ give the price
\[
 A\prod_{i\in I}c_kp_i
\]
for the term with nonempty index set $I$. The $2^k$ factor pays the
product rule; $\mathrm{budC}(k)p_i$ pays the exponential increment when
$p_i\le1$. Its support is $Q\cup\bigcup_{i\in I}X_i$. Since these
supports share their marked site, their covering numbers satisfy
\[
 \operatorname{cov}(Q\cup\bigcup_{i\in I}X_i)
 \le\operatorname{cov}(Q)+\sum_{i\in I}\operatorname{cov}(X_i).
\]
If $\sum_i M_{X_i}(p_i)\le\delta$ and $c_k\delta\le1$, positivity
and the exact finite-product expansion therefore give
\begin{align*}
 \sum_{\varnothing\ne I}M_{Q\cup\cup_{i\in I}X_i}
       (A\prod_{i\in I}c_kp_i)
 &\le M_Q(A)\left[\prod_i(1+c_kM_{X_i}(p_i))-1\right]\\
 &\le M_Q(A)(e^{c_k\delta}-1)
 \le2c_k\delta M_Q(A).
\end{align*}
This also proves regularity term by term: the increments and the local
factor are regular, and the finite product is regular on the displayed
union. No unweighted cardinality estimate is substituted for this sum.

For the actual shifted activity family, the individual prices are the
sum of the exterior-path increment price and the local regular-activity
increment price. Summing them gives
\[
 \delta\le\mathrm{beamAmt}+\mathrm{vat}\,
                  (2r+1)^4\mathrm{hue}(t)\le t^{-109/100}.
\]
Here $r$ is the retained truncation radius. The exterior contribution
contains $e^{-\sigma R/4}$, with fixed $\sigma>0$ and
$R=1000\lceil\log^2t\rceil$, times fixed powers of $t$ and logarithms,
so is smaller than every negative power of $t$.
The regular contribution is bounded by $t^{-11/10+o(1)}$: the radius factors
are subpolynomial under the retained tower bound, and its activity
receipt obeys $\mathrm{hue}(t)\le t^{-11/10+o(1)}$. Thus both terms are
bounded together by $t^{-109/100}$ after increasing the stage threshold.
The local shifted factor, after the fixed derivative composition, has
price $A\le t^{-1/2+o(1)}$. Its support has at most $\mathrm{TileNo}$
covering blocks. These facts give the per-link, per-difference-node
nonempty price
\[
 S(t)\le2c_k t^{-109/100}
                e^{\alpha\mathrm{TileNo}}t^{-1/2+o(1)}.
\]
There are $t^{o(1)}$ marked links, the fixed finite-difference weights
have total absolute value $\sum_{j<m}|w_j|$, and the inverse difference
step is bounded by $t^{1/2+1/24+o(1)}$. Consequently the \emph{actual} support sum,
including the difference weights and normalization, is bounded by
\[
 t^{-109/100-1/2+1/2+1/24+o(1)}
   =t^{-629/600+o(1)}\le t^{-101/100}
   \le\theta^2t^{-1/100}.
\]
The last inequality is the retained threshold bound $t^{-1}\le\theta^2$.
The strict gap between $629/600$ and $101/100$ absorbs the fixed stage
constants. This proves \eqref{eq:trace-charge-tail-price} for the finite
activity family, uniformly in its finite truncations, with no additional
factor $H^2$.

\subsubsection{From the constructed probes to healthy derivative prices}
\label{subsec:trace-charge-healthy-prices}

The local stress, its divergence, and the elementary correction probes
are the entries and first spatial differences of the same localized
polynomial field: the localized stress (\texttt{Meld}), the two fixed
quadratic probes (\texttt{Plint}), the trace current (\texttt{pluck})
and the antisymmetric current (\texttt{pluckD}). This includes their
tensor, single-derivative and divergence entries through grade five. Their local analytic construction gives derivative
bounds $A\le H^{1/50}$ on translated link-coordinate balls of the fixed
collar radius $B\ge1$. Each probe is gauge invariant, local, and has
zero constant and linear Taylor coefficients at the flat field. The
last property must be retained in passing to a healthy chart: an
arbitrary bounded local observable would not have the factor $\theta^2$.

The healthy axial chart has size
$s=\theta(1+r)^p$, with $r$ the support radius and $p$ the fixed stage
power. On this chart use the weighted derivative norm in which a
$j$th derivative is multiplied by $s^j$. Taylor's formula twice gives
$|f(su)|\le C_B A s^2$, and once gives
$s\|Df(su)\|\le C_B A s^2$. For $j\ge2$, the original bound gives
$s^j\|D^jf(su)\|\le A s^2$ since $s\le1$. Thus the entire weighted
norm through the retained order is at most $C_B A s^2$. Pullback through
the healthy collar chart, by the finite chain rule, costs at most
$k!(1+B)^k$. It follows that the healthy derivative price is
\[
 k!(1+B)^k C_B A\theta^2(1+r)^{2p}.
\]
For $r=O(\log^2t)$ the last polynomial is $H^{o(1)}$. Since
$1/50<1/8$, the displayed expression is at most
$\theta^2H^{1/8}$. At the origin of each translated coordinate chart,
the original zeroth-order bound also gives the global estimate
$|f(U)|\le A\le H^{1/50}$. This proves the two distinct prices used in
Subsubsection~\ref{subsec:trace-reference-atomic-certificate} from the local
probe bounds, on the original full observable.

For the distinguished divergence subtract its value at the flat field.
Its first jet vanishes by global conjugation invariance and simplicity
of the gauge Lie algebra: restricting that derivative to any one link
gives an adjoint-invariant linear functional, which vanishes on all
commutators, hence on the entire simple Lie algebra. The link-coordinate
derivatives therefore all vanish. Subtracting the constant supplies the
zero value, and does not alter this first-derivative argument.
Its derivative price before the subtraction is $H^{1/200}$, and the
subtracted constant has absolute value at most this same number. Hence
$A\le2H^{1/200}$ suffices, and the same argument again gives the healthy
price $\theta^2H^{1/8}$. For the small loop counterterm the corresponding
analytic price is $\mathrm{calorie}(t)$; its decay is faster than the
fixed $H^{-2}$ price after the same polynomial chart factors. The
healthy counterterm price is consequently $\theta^2H^{-2}$. Finally,
if a residual or Taylor error has already been bounded in the axial
chart by $\theta^2t^{-1/100}$, the finite pullback factor gives
\[
 k!(1+B)^k\theta^2t^{-1/100}\le\theta^2H^{-2}.
\]
All constants here belong to the selected finite stage and are paid by
the same prior stage-delay envelope as the atomic construction.

\subsubsection{Size of the moment counterterms}
\label{subsec:trace-charge-counterterm-size}

The zero and first moments used to correct the finite loop family are
small for a reason independent of their exact cancellation. The infinite loop family of the selected retained path has exactly
zero zeroth and first moments, by its path law. Restricting that same
family to the radius $R$ loses at most
\[
 e^{-\sigma R/4}\operatorname{momentPrice}(1,\sigma/2)
                \mathrm{humble}(t)
\]
in the combined zeroth- and first-moment norm: insert the factor distance
into the absolutely
summable loop coefficients and spend half of the exponential decay to
pay it. For each fixed $t$, finite truncation of that summable family
can be chosen so that its additional moment error is at most
$t^{-100}$. The combined moment error is therefore bounded by
$t^{-98}$ for all sufficiently large $t$, uniformly after this finite
truncation threshold.

The explicit correction in
Subsubsection~\ref{subsec:trace-charge-loop-cancellation} depends linearly on
these finitely many moments. Its words have length at most ten and
fixed offsets $0,e_a$. Differentiating any such word a fixed number of
times costs a fixed constant; translating it over the support costs at
most $C(1+\mathrm{puff}(t))^4$. Hence its analytic derivative price is
at most $t^{-95}$, after enlarging the stage threshold. Applying the
localized vector field in the directional derivative adds only the
fixed product-rule constant and its subpolynomial derivative price.
The resulting counterterm $d_b$ has analytic price at most $t^{-94}$,
and in particular at most $\mathrm{calorie}(t)=t^{-30}$.
Subsubsection~\ref{subsec:trace-charge-healthy-prices} now gives its healthy
price $\theta^2H^{-2}$. This establishes smallness of the same finite
counterterms which enforce the exact moments; no second family is
chosen for the size bound.

\subsubsection{The fifth-order residual and finite-difference remainder}
\label{subsec:trace-charge-taylor-remainders}

Let $s=\theta(1+r)^p$ be the same axial-chart scale. The action residual
is the actual difference
\[
 K_b=\operatorname{div}M_b-D F_{\rm main}[Y_b^R],
 \qquad f_b^0=\beta K_b.
\]
Here the small four-jet input follows from the inherited action identity,
not from locality alone. Let $J,Y,F$ be the original stress, current and
full main interaction, and let $F_L$ be the same main interaction
restricted to the retained finite family and radius. In this paragraph
$\operatorname{div}M$ is the spatial finite-difference divergence of
the stress, as opposed to the field-space divergence of $Y$ in the
current-price calculation. Then exactly
\[
 \operatorname{div}M-DF_L[Y^R]
 = (\operatorname{div}J-DF[Y])
    +\operatorname{div}(M-J)
    +DF[Y-Y^R]+D(F-F_L)[Y^R].
\]
The first summand has weighted four-jet price at most $t^{-100}$ by
the inherited action equation of the selected field. The second has
price at most $t^{-70}$ times a fixed polynomial in $a^{-1}$ by the
localization estimate. For the third, the weighted product rule pairs
the localized-current tail $t^{-70}$ with the full action's link
derivative price through four further field derivatives (thus through
five derivatives of the scalar action), bounded by a fixed power of
$H$ and $a^{-1}$. The retained action family has exactly this bound;
its four-derivative bound for the scalar action alone would not suffice. Finally absolute
summability of the action family permits its finite truncation error to
be at most $t^{-100}$, and the radius restriction costs at most
$C\operatorname{momentPrice}(P,\sigma/2)e^{-\sigma R/4}$ times the
same fixed support price. Multiplying this by the original-current
price plus its $t^{-70}$ tail bounds the last summand. All spatial
contractions here are exponentially weighted link sums; their moment
losses are polynomials in $a^{-1}$, independent of total volume.
The sum is $t^{-70+o(1)}+t^{-100+o(1)}$, hence at most $t^{-30}$
eventually. This proves the claimed origin four-jet bound for the exact
$K_b$ above. Meanwhile the local analytic construction
bounds all derivatives through the required fixed order by $C H^{q_0}$
for fixed $C,q_0$
on the radius-$Bs$ ball. Both statements concern the same finite main
interaction: its finite truncation is chosen after the local four-jet
comparison. For $0\le j\le4$ and sufficiently large $t$,
\[
 t^{-30}\le C H^{q_0}s^{5-j},\qquad t^{-1}\le s\le1.
\]
Taylor's formula for $D^jK_b(su)$ through degree $4-j$, with the fifth
and higher derivatives bounded on the line segment, therefore bounds
each scaled derivative by $C_B C H^{q_0}s^5$. For $j\ge5$, the same
bound follows directly from $s^j\le s^5$. The finite Taylor constants
are the fixed fifth-order interpolation constant $C_B$. Multiplying by
$|\beta|\le4t$ gives
\[
 \|\beta K_b(s\,\cdot)\|_{C^k(B)}
 \le4t\,C_B C H^{q_0}s^5
 \le\theta^2t^{-1/100}.
\]
Indeed, after division by $\theta^2$, the left side is
$t^{-1/2}$ times a fixed power of $H$ and $1+r$, hence is
$t^{-1/2+o(1)}$. This is the actual fifth-order residual estimate; the
quadratic healthy-price argument alone would not prove it.

For the finite-difference remainder let $h=\theta t^{-1/24}$ be the
chosen step and $m\ge48$ its order. At each marked link and color,
write the shifted integrand as the product of the local current factor
and the exponential action increment. The latter increment is zero at
zero shift. The derivative bound for that increment is
$Z(t)=2\,\mathrm{shear}(t)C_Bs^2\le H^c$. The segment has length at
most $\kappa_m h/s$, where $\kappa_m$ depends only on the fixed
finite-difference stencil. Thus
$Z(t)\kappa_m h/s\le H^c\kappa_m t^{-1/24}\le1$ eventually.
Integrating its first derivative along this segment bounds the increment
itself by one. The scalar exponential has bounded outer derivatives on
$[-1,1]$; after composition its derivatives through order $k+m+1$ are
bounded by fixed powers of $H$, not by an absolute constant. The chain
rule and the product rule therefore bound the full integrand by
\[
 C(t)s,\qquad C(t)\le H^c
\]
for some fixed $c$, after including the local-current price and the
finite coordinate-chart factors. For clarity, the action-dependent
composition factor contains $\mathrm{shear}(t)s^2$; since
$\mathrm{shear}(t)\le5tH^{1/50}$ and
$s^2=t^{-1}H^2(1+r)^{2p}$, this factor is bounded by a fixed power of
$H$. All further powers are fixed by $k+m+1$.

The finite-difference weights cancel the Taylor coefficients through
order $m$. Taylor's integral remainder, differentiated up to order $k$,
then bounds the error at each marked link and color by
$C(t)s(h/s)^{m+1}$. Summing over the actual finite support multiplies
$C(t)$ by at most $O((1+r)^4)\dim\mathfrak g$, still a fixed power of
$H$. Since $h/s\le t^{-1/24}$,
\[
 C(t)s(h/s)^{m+1}
   = C(t)h(h/s)^m
   \le t^{-2+o(1)}h
   \le\theta^2t^{-1/100}h.
\]
The observable $f_b^3$ divides this Taylor error by $h$. That division
cancels the displayed factor $h$ exactly, leaving the same axial price
$\theta^2t^{-1/100}$ as for $f_b^0$. The finite collar pullback established
above therefore gives the healthy price $\theta^2H^{-2}$ for both errors.
The finite-difference integrand estimate also proves the shifted local
factor price $t^{-1/2+o(1)}$ used in
Subsubsection~\ref{subsec:trace-charge-support-tail}: it is $s$ times fixed
powers of $H$, before taking the Taylor remainder or dividing by $h$.
